\subsubsection{可逆性阈值：\texorpdfstring{$\Phi_m$}{Phi\_m} 的共轭与纤维坍缩}\label{subsec:Phi_m-conjugacy}

上一节已将 $Y_m$ 组织为标记图 $G_m$ 的像（定理 \ref{thm:Phi_m-sofic-graph}），并证明熵率不降（推论 \ref{cor:Phi_m-entropy-no-drop}）。
本节给出可逆性与纤维结构的确定结论：在 Zeckendorf 折叠 $\Fold_m$ 的口径下，除去 $m=2$ 的两点退化外，$\Phi_m$ 在 $m\ge 3$ 时为共轭。

\paragraph{可终止的局部判定口径（子集滤波）}
沿定义 \ref{def:G_m}，对任意候选后缀集合 $S\subseteq V_m$ 与输出符号 $a\in X_m$，定义一步更新
$$
\delta_m(S,a):=\{v'\in V_m:\exists\,e:v\to v'\ \text{s.t. }v\in S,\ \ell_m(e)=a\}.
$$
从初始集合 $S_0:=V_m$ 出发，对任意输出块 $a_0a_1\cdots a_{t-1}$ 递推
$$
S_{j+1}:=\delta_m(S_j,a_j)\qquad(0\le j<t).
$$
该递推可视为“仅由输出块驱动”的有限状态滤波器：$S_j$ 记录读入前 $j$ 个输出符号后，内部后缀（顶点）仍可能取到的候选集。
当 $S_t$ 为单点集时，意味着该长度-$t$ 输出块已将内部后缀同步到唯一值。

\begin{theorem}[可逆性阈值与有限窗逆码]\label{thm:Phi_m-conjugacy-threshold}
令 $\Phi_m$ 为定义 \ref{def:Phi_m} 的滑动块码，$Y_m$ 为定义 \ref{def:Y_m} 的像子移位。
则：
\begin{enumerate}
  \item 当 $m\ge 3$ 时，$\Phi_m:\{0,1\}^{\ZZ}\to Y_m$ 为拓扑共轭；更具体地，存在一个滑动块码
  $$
  \Psi_m:Y_m\to\{0,1\}^{\ZZ},
  $$
  其\textbf{记忆}为 $m-1$、\textbf{预期性}为 $0$，并满足
  $$
  \Psi_m\circ\Phi_m=\mathrm{id}_{\{0,1\}^{\ZZ}},\qquad
  \Phi_m\circ\Psi_m=\mathrm{id}_{Y_m}.
  $$
  \item 当 $m=2$ 时，$\Phi_2$ 恰好在两条常值序列 $0^{\ZZ}$ 与 $1^{\ZZ}$ 上出现 $2$-对-$1$ 退化，其余点一一。
\end{enumerate}
\end{theorem}

\begin{proof}
证明思路建立在“有限可判定”的滤波器上。

\paragraph{(A) 局部同步：长度-$m$ 块把候选集压到单点}
对固定 $m$，考虑上述滤波更新 $S_{j+1}=\delta_m(S_j,a_j)$，并记
$$
R_m:=\{S_m:\exists\,a_0\cdots a_{m-1}\in\mathcal{L}_m(Y_m)\ \text{s.t. 递推得到 }S_m\}.
$$
若 $R_m$ 中每个元素都是单点集，则任意长度-$m$ 的输出块都把内部后缀同步到唯一顶点。
该命题对每个固定 $m$ 是一个\textbf{完全有限}的判定问题：状态空间为 $V_m$ 的幂集可达部分，转移由 $G_m$ 给出。
脚本 \texttt{scripts/exp\_phi\_m\_conjugacy\_threshold.py} 对 $m\in[2,12]$ 的范围给出可复算证书，输出汇总见表 \ref{tab:phi_m_conjugacy_threshold}。
\begin{table}[H]
\centering
\scriptsize
\setlength{\tabcolsep}{6pt}
\caption{Conjugacy threshold certificate for the sliding-block factor $\Phi_m$. We build the subset-filter automaton from the labeled graph $G_m$ and compute $R_m$, the set of reachable filter states after exactly $m$ output symbols. If all states in $R_m$ are singletons, then each length-$m$ output word has a unique length-$(2m-1)$ preimage block, yielding a finite-window inverse (memory $m-1$). We also sanity-check periodic-point counts up to a small $n$.}
\label{tab:phi_m_conjugacy_threshold}
\begin{tabular}{r r r r r l r}
\toprule
$m$ & det states & det edges & $|R_m|$ & non-singleton ends & periodic $\#\mathrm{Fix}_n$ & checked $n$\\
\midrule
2 & 3 & 7 & 3 & 1 & $2^n-1$ & 8\\
3 & 7 & 20 & 4 & 0 & $2^n$ & 8\\
4 & 16 & 50 & 8 & 0 & $2^n$ & 8\\
5 & 32 & 95 & 16 & 0 & $2^n$ & 8\\
6 & 69 & 204 & 32 & 0 & $2^n$ & 8\\
7 & 142 & 404 & 64 & 0 & $2^n$ & 8\\
8 & 296 & 827 & 128 & 0 & $2^n$ & 8\\
9 & 607 & 1656 & 256 & 0 & $2^n$ & 8\\
10 & 1244 & 3342 & 512 & 0 & $2^n$ & 8\\
11 & 2532 & 6703 & 1024 & 0 & $2^n$ & 8\\
12 & 5141 & 13464 & 2048 & 0 & $2^n$ & 8\\
\bottomrule
\end{tabular}
\end{table}

其中 $m\ge 3$ 时 $R_m$ 的“non-singleton ends $(m)$”为 $0$，而 $m=2$ 时在任何步数都出现唯一一个非单点端态，对应常值输出块所导致的两点退化。
同时表中 “non-singleton ends $(m-1)$” 在 $m\ge 3$ 时严格为正，这给出同步延迟的尖锐性见证：一般需要完整的长度-$m$ 输出块才能把候选集压到单点（对应逆码记忆为 $m-1$）。

\paragraph{(B) 从单点端态到唯一 lift：利用双向解析性}
由定理 \ref{thm:pom-right-resolving} 与定理 \ref{thm:pom-left-resolving}，标记图 $G_m$ 在 Zeckendorf 折叠下同时右解析与左解析：同一顶点出发（或到达）的一步扩展可由标记唯一判别。
因此一旦长度-$m$ 输出块把端态同步为单点 $\{v'\}$，则该块在 $G_m$ 上的长度-$m$ lift 路径不仅存在，而且由左解析性可从终点 $v'$ 逐步唯一回溯得到；从而该输出块对应唯一的长度-$(2m-1)$ 输入块。

\paragraph{(C) 构造逆码 \texorpdfstring{$\Psi_m$}{Psi\_m}}
当 $m\ge 3$ 时，由 (A)(B) 可知：对任意允许的长度-$m$ 输出块
$$
a_{t-m+1}\cdots a_t\in\mathcal{L}_m(Y_m),
$$
存在唯一的 $G_m$-路径
$$
e_{t-m+1},\dots,e_t\in E_m
$$
使得其逐边标记恰为该输出块。
令 $e_{t-m+1}=(v,b)$（其中 $b\in\{0,1\}$ 是该边对应的追加比特），定义
$$
(\Psi_m(y))_t:=b,
$$
其中 $y\in Y_m$ 且 $a_j=y_j$。
该定义仅依赖 $y_{t-m+1},\dots,y_t$，故 $\Psi_m$ 的记忆为 $m-1$、预期性为 $0$。
由构造立得 $\Psi_m\circ\Phi_m=\mathrm{id}$，并因 $\Phi_m$ 满射到 $Y_m$ 得 $\Phi_m\circ\Psi_m=\mathrm{id}_{Y_m}$。

\paragraph{(D) $m=2$ 的两点退化}
脚本证书给出：唯一的非单点端态对应常值输出块，其恰来自两条输入常值序列 $0^{\ZZ},1^{\ZZ}$；
此外对任意其它输出点滤波在有限步内同步到单点，从而 lift 唯一。
综上得到 $m=2$ 的“仅两点 $2$-对-$1$”退化结论。
\end{proof}

\begin{corollary}[\texorpdfstring{$\zeta_{Y_m}$}{zetaYm} 的封口与“折叠因子不产出非平凡 prime shadow”]\label{cor:Phi_m-zeta-closed}
对 $m\ge 3$，由于 $Y_m$ 与二元全移位共轭（定理 \ref{thm:Phi_m-conjugacy-threshold}），其周期点计数满足
$$
\#\mathrm{Fix}(\sigma^n|_{Y_m})=2^n\qquad(n\ge 1),
$$
从而内禀 Artin--Mazur \(\zeta\)-函数为
$$
\zeta_{Y_m}(z)=\frac{1}{1-2z}.
$$
而 $m=2$ 时有
$$
\#\mathrm{Fix}(\sigma^n|_{Y_2})=2^n-1,\qquad
\zeta_{Y_2}(z)=\frac{1-z}{1-2z}.
$$
因此当 $m\ge 3$ 时，$Y_m$ 自身的 prime-orbit/\(\zeta\) 指纹退化为全移位骨架；本文中出现的非平凡 prime shadow 必然来自“核 + 势/权重 + 扭曲通道”（而非单纯来自折叠因子 $Y_m$）。
\end{corollary}
